%============================================================================
% BASALTHQ PROPOSAL — OFFICE OF THE CTO
% BasaltDelivers: Two-Touchpoint Delivery Platform on Surge
%============================================================================
\documentclass[11pt, letterpaper]{article}
\usepackage{basalthq-brand}
\usepackage{tocloft}
\usepackage{needspace}
\renewcommand{\cftsecleader}{\cftdotfill{\cftdotsep}}

% ============================================================================
% DOCUMENT VARIABLES
% ============================================================================
\setauthor{Krishna Patel}{Chief Technical Officer \& Founder}
\setdepartment{Office of the CTO}{DeptCTO}
\setclassification{Confidential}

\newcommand{\proposalTitle}{BasaltDelivers}
\newcommand{\proposalSubtitle}{A two-touchpoint delivery ecosystem on BasaltSurge --- web, Android, and iOS}
\newcommand{\proposalTo}{Eric Turner, Chief Executive Officer}
\newcommand{\proposalVersion}{1.2}
\newcommand{\proposalBudget}{\$0 (Internal Engineering)}

% ============================================================================
\begin{document}
% ============================================================================

\setupheaderfooter{Proposal --- \proposalTitle}
\raggedbottom
\widowpenalty=10000
\clubpenalty=10000

% ============================================================================
% COVER PAGE
% ============================================================================
\begin{titlepage}
\begin{center}
    \vspace*{3cm}

    {\fontsize{14}{18}\selectfont\color{BasaltText}\textsc{BasaltHQ Inc.}}

    \vspace{0.5em}
    {\color{\docDeptColor}\rule{0.4\linewidth}{2.5pt}}

    \vspace{2em}
    {\fontsize{12}{16}\selectfont\color{BasaltText}\textsc{Proposal}}

    \vspace{0.5em}
    {\fontsize{32}{40}\selectfont\bfseries\color{BasaltDark} \proposalTitle}

    \vspace{0.8em}
    {\fontsize{13}{17}\selectfont\color{BasaltText}\itshape \proposalSubtitle}

    \vspace{4em}

    \renewcommand{\arraystretch}{1.4}
    \begin{tabular}{rl}
        {\small\textsc{Prepared by}} & \docAuthorName, \docAuthorTitle \\
        {\small\textsc{Prepared for}} & \proposalTo \\
        {\small\textsc{Department}} & \docDepartment \\
        {\small\textsc{Date}} & \today \\
        {\small\textsc{Version}} & \proposalVersion \\
        {\small\textsc{Budget}} & \proposalBudget \\
        {\small\textsc{Classification}} & \docClassification \\
    \end{tabular}
    \renewcommand{\arraystretch}{1.0}

    \vfill
    {\color{\docDeptColor}\rule{0.4\linewidth}{2.5pt}}

    \vspace{1em}
    {\small\color{BasaltText}\textsc{BasaltHQ Inc. \quad$\cdot$\quad Murrieta, California}}
\end{center}
\end{titlepage}

\tableofcontents
\newpage

% ============================================================================
% PROPOSAL BODY
% ============================================================================

\section{Executive Summary}

This proposal introduces \textbf{BasaltDelivers}, a delivery-as-a-service product built entirely on the existing BasaltSurge protocol. The product is deployed as two \emph{touchpoints} within the Surge ecosystem, each available on web, Android, and iOS:

\begin{enumerate}
    \item \textbf{BasaltDelivers} --- the consumer-facing storefront where customers browse merchant offerings and place delivery orders.
    \item \textbf{BasaltDrive} --- the driver-facing operations hub where independent delivery drivers sign up, receive dispatch requests, navigate to delivery locations, and confirm completed deliveries.
\end{enumerate}

Both touchpoints are first built as route groups within the existing \texttt{portalpay-official} Next.js application, then packaged as native Android and iOS applications using \textbf{Capacitor} (by Ionic). This ``write once, deploy everywhere'' strategy gives us App Store and Play Store presence---with full access to native device APIs (GPS, push notifications, camera)---while sharing 100\% of the business logic with the web touchpoints.

BasaltDelivers and BasaltDrive are monetized as \textbf{add-on subscription services} at \textbf{\$50/month per merchant}, billed via Stripe. This replaces any per-delivery commission model---merchants retain 100\% of their delivery fees and tips, and compensate their drivers independently. The standard BasaltSurge platform fee (moving to 1\%) continues to apply to the underlying transaction as usual.

Dedicated vanity domains (\texttt{basaltdelivers.com} and \texttt{basaltdrive.com}, pending finalization) will forward to the Surge-hosted web versions. Because the entire product is built on existing Surge infrastructure and AI-assisted code generation tools (Antigravity, Claude Code) compress our development velocity by an estimated 5--8$\times$, the total time-to-market is approximately \textbf{4 weeks}.

\begin{actionbox}
\textbf{Decision Requested:} Approve the creation of the two Surge touchpoints with Capacitor-based native app packaging, Stripe subscription billing integration, domain procurement, and Apple Developer Program enrollment (\$99/year). All engineering is AI-assisted internal work at zero labor cost.
\end{actionbox}

\section{Problem Statement}

\subsection{Current State}
BasaltSurge currently serves merchants through point-of-sale terminals, handhelds, kiosks, and online shop portals. While merchants can manage full-service dine-in and takeout operations, the platform does not yet offer a first-party delivery channel. Merchants who wish to offer delivery must rely on third-party aggregators (DoorDash, Uber Eats, Grubhub) that extract 15--30\% commission per order, eroding margins and fragmenting the customer relationship.

On the driver side, no mechanism exists within Surge for independent drivers to discover, apply to, and fulfill delivery orders for our onboarded merchants.

\subsection{Impact}
\begin{itemize}
    \item \textbf{Revenue leakage:} Every order fulfilled through a third-party aggregator costs the merchant 15--30\% in commission---revenue that could be retained under the Surge 1\% model.
    \item \textbf{Brand dilution:} Merchants lose direct customer relationships when orders flow through aggregator apps.
    \item \textbf{Platform gap:} Without a native delivery channel, BasaltSurge is incomplete as an end-to-end commerce protocol, limiting our competitive positioning against full-stack alternatives.
    \item \textbf{Driver ecosystem void:} There is no pathway for gig-economy drivers to participate in our ecosystem, representing a missed network-effect opportunity.
    \item \textbf{Subscription revenue opportunity:} A \$50/month add-on service creates predictable, recurring SaaS revenue independent of transaction volume---a fundamentally more stable revenue stream than per-delivery commissions.
\end{itemize}

\section{Proposed Solution}

\subsection{Approach}
We will implement BasaltDelivers as two dedicated touchpoint routes within the existing \texttt{portalpay-official} Next.js application, leveraging the multi-tenant brand resolution system already in production. Each touchpoint inherits the full Surge infrastructure stack:

\begin{itemize}
    \item \textbf{Inventory API} for menu and catalog browsing
    \item \textbf{Orders API} for unified order lifecycle management
    \item \textbf{Receipts API} for financial records and guest-facing portals
    \item \textbf{BPS Fee Engine} for the 1\% platform fee on underlying transactions
    \item \textbf{Thirdweb Authentication} for identity and wallet-based sessions
    \item \textbf{Stripe Billing} for \$50/month merchant subscription management
    \item \textbf{WebSocket Infrastructure} for real-time order status and driver location streaming
    \item \textbf{Cloudflare DNS \& Plesk Hosting} for domain forwarding and SSL termination
\end{itemize}

\subsubsection{Domain Architecture}
Two vanity domains will be procured and configured:

\begin{center}
\begin{tabularx}{\textwidth}{@{}l l X@{}}
\toprule
\textbf{Domain} & \textbf{Forwards To} & \textbf{Purpose} \\
\midrule
\texttt{basaltdelivers.com} & \texttt{surge.basalthq.com/delivers} & Consumer delivery storefront \\
\texttt{basaltdrive.com} & \texttt{surge.basalthq.com/drive} & Driver operations hub \\
\bottomrule
\end{tabularx}
\end{center}

Both domains will CNAME to \texttt{surge.basalthq.com} and leverage the existing Cloudflare SSL-for-SaaS pipeline for automatic certificate provisioning. The brand resolution system will detect the incoming hostname and apply the appropriate BasaltDelivers or BasaltDrive branding context.

Critically, this requires only a \textbf{slight modification} to our existing custom domain system. Surge already supports merchant-to-touchpoint domain mapping via the Cloudflare DNS manager (\texttt{cloudflare/manager.ts}) and Plesk nginx routing. Today, merchants CNAME their own domains to \texttt{surge.basalthq.com}, and the platform resolves the hostname to the correct brand context. For BasaltDelivers, we simply register \texttt{basaltdelivers.com} and \texttt{basaltdrive.com} as platform-level domain entries in the same resolution table---rather than merchant-scoped entries---and map them to the \texttt{(delivers)} and \texttt{(drive)} route groups respectively. No new DNS infrastructure or SSL pipeline is needed.

\subsection{Key Deliverables}

\begin{enumerate}
    \item \textbf{BasaltDelivers Consumer Touchpoint:} A responsive web application where consumers can browse participating merchants, view menus and item details, compose delivery orders with modifier support (including merchant-configured delivery fees), and complete checkout via the embedded Surge payment portal. Real-time order tracking with driver location streaming is included.
    \item \textbf{BasaltDrive Driver Touchpoint:} A responsive web application where drivers can create profiles, discover merchants within their geolocation radius, submit applications to join merchant delivery teams, receive real-time delivery dispatch notifications via WebSocket, access turn-by-turn navigation to the delivery address, and confirm delivery completion with photo proof.
    \item \textbf{Admin Team Panel --- Driver Role:} A new ``Drivers'' tab within the existing merchant Team panel in the admin module, adding a \textbf{Driver} role to the existing team role system. The tab includes three subtabs:
    \begin{itemize}
        \item \textbf{Active} --- Currently approved and operational drivers with real-time status, vehicle info, and performance metrics.
        \item \textbf{Requested} --- Incoming driver applications (filtered by the merchant's configured delivery radius) with one-click approve/deny actions.
        \item \textbf{Denied} --- Previously rejected applicants with the option to reconsider.
    \end{itemize}
    Merchants can also configure delivery fee amounts, attire requirements, and delivery scripts from this panel.
    \item \textbf{Stripe Subscription Billing:} Integration of Stripe Checkout and Customer Portal for \$50/month BasaltDelivers add-on billing, with automated feature gating (delivery touchpoints activate only for merchants with an active subscription).
    \item \textbf{Domain Forwarding Infrastructure:} DNS configuration for \texttt{basaltdelivers.com} and \texttt{basaltdrive.com} with SSL termination and brand-context injection via the existing Cloudflare + Plesk hosting stack.
\end{enumerate}

\subsection{Alternatives Considered}

\begin{center}
\begin{tabularx}{\textwidth}{@{}l X X@{}}
\toprule
\textbf{Option} & \textbf{Pros} & \textbf{Cons} \\
\midrule
Option A --- Surge + Capacitor (Recommended) & Single codebase for web, Android, and iOS; full native API access (GPS, push, camera); ships to App Store and Play Store; zero logic duplication; AI-accelerated development & Requires Apple Developer enrollment (\$99/yr); Capacitor plugin ecosystem for edge-case native APIs \\
Option B --- Surge Web Only & Simplest implementation; no app store overhead & No push notifications; degraded GPS on mobile browsers; no discoverability in app stores \\
Option C --- React Native / Expo & Near-native performance; cross-platform from one codebase & Requires rewriting all UI from React/Next.js to React Native primitives; separate build toolchain; doubles maintenance surface \\
Option D --- Native Swift + Kotlin & Best possible UX and performance & Two completely separate codebases; requires macOS for iOS builds; 3--6 month timeline even with AI tools \\
Option E --- Do Nothing & No engineering cost & Merchants continue losing 15--30\% to aggregators; platform remains incomplete; driver network opportunity forfeited \\
\bottomrule
\end{tabularx}
\end{center}

\begin{insightbox}
\textbf{Why Capacitor over React Native:} Our touchpoints are already Next.js web applications. Capacitor wraps the \emph{existing} web app in a native shell and provides a bridge to native device APIs via plugins---no UI rewrite required. The same \texttt{npm run build} output that serves the web also powers the native apps. Capacitor's plugin ecosystem covers all critical requirements: \texttt{@capacitor/geolocation} for high-accuracy GPS, \texttt{@capacitor/push-notifications} for real-time dispatch alerts, \texttt{@capacitor/camera} for proof-of-delivery photos, and \texttt{@capacitor/haptics} for tactile feedback on driver actions.
\end{insightbox}

\section{Scope \& Timeline}

\begin{warningbox}
\textbf{AI-Accelerated Development:} This timeline assumes the use of state-of-the-art agentic code generation tools (Antigravity, Claude Code) which provide 5--8$\times$ velocity gains on scaffolding, API integration, UI component generation, and boilerplate-heavy tasks like Capacitor plugin wiring. A single engineer equipped with these tools can produce output equivalent to a 4--6 person team working conventionally. The phases below reflect this accelerated cadence.
\end{warningbox}

\begin{center}
\begin{tabularx}{\textwidth}{@{}l l l X@{}}
\toprule
\textbf{Phase} & \textbf{Duration} & \textbf{Milestone} & \textbf{Deliverables} \\
\midrule
Phase 1 & Days 1--3 & Infrastructure \& APIs & Database schema for \texttt{drivers}, \texttt{empanelments}, and \texttt{delivery\_orders} in the \texttt{skynetpod} partition; full REST API surface for driver registration, geolocation-based discovery, empanelment lifecycle, and delivery order management; Stripe subscription product and webhook configuration; domain CNAME setup; Capacitor project initialization for both apps \\
Phase 2 & Days 4--8 & Consumer Touchpoint & BasaltDelivers storefront UI with merchant discovery, menu browsing, cart with modifier and delivery-fee support, Surge payment portal checkout, and real-time order tracking; Capacitor build verified on Android emulator \\
Phase 3 & Days 9--14 & Driver Touchpoint & BasaltDrive dashboard with geolocation-radius merchant discovery, team application flow, WebSocket-based delivery dispatch, native GPS tracking via \texttt{@capacitor/geolocation}, map-based navigation with deep-link to native Maps, proof-of-delivery camera capture, and delivery confirmation; Capacitor build verified on iOS simulator \\
Phase 4 & Days 15--19 & Team Panel \& Native Polish & Admin Team panel ``Drivers'' tab with Active/Requested/Denied subtabs; delivery fee and attire/script configuration; Stripe subscription gating and Customer Portal integration; push notification wiring via \texttt{@capacitor/push-notifications}; native splash screens and app icons for both platforms \\
Phase 5 & Days 20--25 & QA, Submission \& Launch & End-to-end QA across web, Android, and iOS; Stripe subscription lifecycle testing; domain SSL verification; App Store and Play Store submissions; production deployment; merchant onboarding documentation \\
\bottomrule
\end{tabularx}
\end{center}

\begin{insightbox}
\textbf{Total elapsed time: $\sim$4 weeks} (25 working days), accounting for App Store review cycles running in parallel with final QA. Compare to the 12--16 week industry standard for a comparable multi-platform delivery product built without AI tooling.
\end{insightbox}

\section{Revenue Model}

BasaltDelivers is monetized as a \textbf{subscription add-on service}, not a per-delivery commission. This creates predictable, recurring SaaS revenue that is independent of transaction volume.

\subsection{Subscription Pricing}
\begin{itemize}
    \item \textbf{Price:} \$50/month per merchant (billed via Stripe).
    \item \textbf{Includes:} Full access to the BasaltDelivers consumer storefront, BasaltDrive driver operations hub, the admin Team panel Driver role, and unlimited delivery orders.
    \item \textbf{Feature Gating:} Delivery touchpoints are only visible and functional for merchants with an active Stripe subscription. Cancellation immediately disables delivery ordering for that merchant while preserving their core Surge POS functionality.
\end{itemize}

\subsection{Payment Flow}
Delivery orders use the \textbf{existing} Surge payment infrastructure with no new split deployment:
\begin{enumerate}
    \item The consumer places a delivery order. The order total includes the menu items plus the \textbf{merchant-configured delivery fee}.
    \item Payment is processed through the standard Surge payment portal. The platform retains its 1\% fee (\texttt{platformBps: 100}) on the full transaction amount, including the delivery fee.
    \item The merchant receives the remainder (99\% of the order total, including the delivery fee and any tip).
    \item The merchant compensates their drivers independently---via cash, Zelle, direct deposit, or any method of their choosing. This is an operational relationship between the merchant and their team, not mediated by the platform.
\end{enumerate}

\begin{insightbox}
\textbf{Why subscription over commission:} A flat \$50/month fee is simpler for merchants to budget, avoids the ``DoorDash anxiety'' of watching per-order commissions erode margins, and gives BasaltHQ predictable MRR. Merchants who process even modest delivery volume (3--4 orders/day at \$30 average) save \$4,000--\$9,000/year vs.\ aggregator commissions while paying only \$600/year to us.
\end{insightbox}

\subsection{Stripe Integration}
Subscription billing is handled entirely by Stripe:
\begin{itemize}
    \item \textbf{Stripe Checkout:} Merchants subscribe via a hosted Stripe Checkout session linked from their admin dashboard.
    \item \textbf{Stripe Customer Portal:} Merchants manage payment methods, view invoices, and cancel subscriptions via the Stripe-hosted Customer Portal.
    \item \textbf{Webhook Enforcement:} A \texttt{/api/stripe/webhook} endpoint listens for subscription lifecycle events (\texttt{created}, \texttt{paid}, \texttt{deleted}) to toggle the \texttt{deliveryEnabled} flag on the merchant's shop configuration document in real time.
\end{itemize}

\needspace{5\baselineskip}
\section{Budget \& Resources}

\begin{center}
\begin{tabularx}{\textwidth}{@{}l r X@{}}
\toprule
\textbf{Line Item} & \textbf{Cost} & \textbf{Notes} \\
\midrule
Engineering (Internal) & \$0 & AI-assisted development by in-house engineering; Antigravity and Claude Code licenses covered under existing subscriptions \\
Domain Registration & \$30/yr & Two \texttt{.com} domains at standard registrar pricing \\
Apple Developer Program & \$99/yr & Required for iOS App Store distribution \\
Google Play Console & \$25 (one-time) & Required for Play Store distribution \\
Stripe Processing & 2.9\% + 30\textcent/txn & Applied to each \$50 subscription charge (\$1.75/merchant/month); net revenue \$48.25/merchant/month \\
Infrastructure (Hosting) & \$0 & Served from existing Plesk VPS; no additional compute required \\
Cloudflare SSL-for-SaaS & \$0 & Covered under existing Cloudflare plan for custom hostnames \\
Database (Cosmos DB) & \$0 & New collections within existing \texttt{skynetpod} partition; no throughput increase needed \\
\midrule
\textbf{Total (Fixed)} & \textbf{\$154 Year 1} & \$129/yr recurring (no repeat Google Play fee) \\
\bottomrule
\end{tabularx}
\end{center}

\subsection{ROI Analysis}

\begin{insightbox}
\textbf{Expected ROI:} Each subscribing merchant generates \textbf{\$579/year} in net subscription revenue (\$48.25/month after Stripe fees). Additionally, all delivery orders continue to generate the standard 1\% platform fee on the transaction amount. For a merchant processing 50 delivery orders/week at \$35 average, the platform earns an additional \$910/year in transaction fees---bringing the total per-merchant yield to approximately \textbf{\$1,489/year}. Meanwhile, the merchant saves \$13,000--\$27,000/year vs.\ aggregator commissions (15--30\% of \$91,000 annual delivery revenue), making the \$600/year subscription a compelling value proposition. At 20 subscribing merchants, the delivery add-on alone generates \$29,780/year in recurring revenue.
\end{insightbox}

\needspace{5\baselineskip}
\section{Technical Architecture}

\subsection{Touchpoint Integration Model}
Both touchpoints are implemented as route groups within the existing \texttt{portalpay-official} Next.js application:

\begin{itemize}
    \item \texttt{src/app/(delivers)/} --- Consumer-facing pages (merchant listing, menu, cart, checkout, tracking)
    \item \texttt{src/app/(drive)/} --- Driver-facing pages (dashboard, empanelment, dispatch, navigation, earnings)
\end{itemize}

The existing \texttt{brand-config.ts} hostname resolution system will be extended to recognize \texttt{basaltdelivers.com} and \texttt{basaltdrive.com} as first-class brand contexts, applying the appropriate layout, color tokens, and navigation structure.

\subsection{Capacitor Native Deployment}
Capacitor (by Ionic) wraps the production web build into native Android and iOS application shells. The architecture is as follows:

\begin{enumerate}
    \item \textbf{Shared Web Core:} The \texttt{next build \&\& next export} output (or static export via \texttt{output: 'export'}) is copied into the Capacitor project's \texttt{dist/} directory. This is the identical artifact served by the web deployment.
    \item \textbf{Native Shell:} Capacitor generates an Xcode project (\texttt{ios/}) and an Android Studio project (\texttt{android/}) that embed the web build in a \texttt{WKWebView} (iOS) or \texttt{WebView} (Android) with full native bridge access.
    \item \textbf{Native Plugin Bridge:} JavaScript code calls native APIs through Capacitor's typed plugin system. No native code is written manually for standard capabilities:
    \begin{itemize}
        \item \texttt{@capacitor/geolocation} --- High-accuracy GPS with background tracking for active deliveries
        \item \texttt{@capacitor/push-notifications} --- FCM (Android) and APNs (iOS) for delivery dispatch alerts
        \item \texttt{@capacitor/camera} --- Photo capture for proof-of-delivery
        \item \texttt{@capacitor/haptics} --- Tactile feedback on driver accept/complete actions
        \item \texttt{@capacitor/app} --- Deep linking from \texttt{basaltdelivers.com} and \texttt{basaltdrive.com} URLs
        \item \texttt{@capacitor/splash-screen} --- Branded launch experience
    \end{itemize}
    \item \textbf{Two Capacitor Projects:} We maintain two separate Capacitor configurations (\texttt{capacitor.config.ts}) to produce distinct app binaries:
    \begin{itemize}
        \item \texttt{com.basalthq.delivers} --- Consumer app (BasaltDelivers)
        \item \texttt{com.basalthq.drive} --- Driver app (BasaltDrive)
    \end{itemize}
    Each configuration sets a different \texttt{server.url} in development mode and a different \texttt{webDir} for production builds, ensuring brand-context isolation at the native level.
\end{enumerate}

\begin{recommendbox}
\textbf{Build Pipeline:} Android builds can be produced on Windows via Android Studio. iOS builds require a macOS environment; we will use a cloud CI service (GitHub Actions with \texttt{macos-latest} runner) for automated iOS archive and App Store Connect uploads, eliminating the need for local macOS hardware.
\end{recommendbox}

\subsection{Data Model Extensions}
New entities will be created within the existing \texttt{skynetpod} database partition:

\begin{itemize}
    \item \textbf{\texttt{drivers}} --- Driver profiles including vehicle type, background check status, real-time GPS coordinates, online/offline state, and home geolocation (used for radius-based merchant discovery).
    \item \textbf{\texttt{empanelments}} --- Many-to-many relationships between drivers and merchants, with lifecycle status (\texttt{requested}, \texttt{active}, \texttt{denied}), attire acknowledgment flags, and script acknowledgment timestamps. These records power the Team panel's Drivers tab subtabs.
    \item \textbf{\texttt{delivery\_orders}} --- Extension of the existing \texttt{orders} entity with delivery-specific fields: driver assignment, merchant-configured delivery fee amount, pickup confirmation timestamp, delivery proof (photo URL), estimated delivery time, and real-time location trail.
    \item \textbf{\texttt{shop\_config} extension} --- New fields: \texttt{deliveryEnabled} (toggled by Stripe webhook), \texttt{deliveryFee} (flat fee added to orders), \texttt{deliveryRadius} (km, for driver discovery filtering), and Stripe identifiers (\texttt{stripeSubscriptionId}, \texttt{stripeCustomerId}).
\end{itemize}

\subsection{Driver Onboarding \& Team Panel Integration}
Driver management is integrated into the \textbf{existing Admin Team panel} rather than a standalone fleet management module. This leverages the role-based team architecture already in production:

\begin{enumerate}
    \item \textbf{New ``Driver'' Role:} A \texttt{driver} role is added to the existing team role enum (alongside \texttt{admin}, \texttt{manager}, \texttt{cashier}, etc.). This role grants access to the BasaltDrive touchpoint for that specific merchant.
    \item \textbf{``Drivers'' Tab:} A new top-level tab appears in the Team panel for merchants with an active BasaltDelivers subscription. It contains three subtabs:
    \begin{itemize}
        \item \textbf{Active} --- Approved drivers currently on the merchant's team. Shows real-time online/offline status, vehicle type, and delivery performance metrics (on-time \%, average delivery time). Merchants can suspend or remove drivers here.
        \item \textbf{Requested} --- Incoming applications from drivers who discovered the merchant via geolocation-radius search in BasaltDrive. Each application shows the driver's profile, vehicle info, distance from the merchant, and a brief message. One-click \textbf{Approve} or \textbf{Deny} actions transition the empanelment record.
        \item \textbf{Denied} --- Previously rejected applicants. Merchants can reconsider and approve from this view if circumstances change.
    \end{itemize}
    \item \textbf{Geolocation-Based Discovery:} When a driver opens BasaltDrive, the app uses \texttt{@capacitor/geolocation} to determine their position and queries the API for merchants within a configurable radius (default 15 km) that have active BasaltDelivers subscriptions. Drivers can then submit applications directly from the discovery view.
    \item \textbf{Delivery Configuration:} From the Drivers tab, merchants also configure their \textbf{delivery fee} (a flat amount added to consumer orders), \textbf{delivery radius} (which controls both consumer eligibility and driver discovery), \textbf{attire requirements}, and \textbf{delivery scripts}.
\end{enumerate}

\subsection{Payment \& Fee Routing}
Delivery transactions flow through the \textbf{existing} BPS fee engine with no new split deployment:
\begin{itemize}
    \item \textbf{Platform fee:} 1\% (\texttt{platformBps: 100}) retained by BasaltSurge on the full order amount (menu items + delivery fee + tip).
    \item \textbf{Delivery fee:} A merchant-configured flat fee appended to the consumer's order total. The merchant receives this fee in full (minus the 1\% platform fee) as part of their standard Surge settlement.
    \item \textbf{Driver compensation:} Handled entirely by the merchant outside the platform. The merchant-driver financial relationship is an internal team arrangement---similar to how a restaurant pays its waitstaff. No new \texttt{splitConfig} or on-chain split contract is required.
    \item \textbf{Tips:} Consumer tips are included in the order total, settled to the merchant, and forwarded to drivers at the merchant's discretion.
\end{itemize}

\begin{recommendbox}
\textbf{Architectural Simplification:} By keeping driver compensation off-platform, we avoid the complexity of a new split deployment, additional \texttt{splitConfig} documents, and the associated synchronization traps documented in the Fee Infrastructure KI. The merchant's existing Surge payment flow is completely unchanged---delivery orders are simply regular orders with an additional line item for the delivery fee.
\end{recommendbox}

\subsection{Real-Time Infrastructure}
The existing WebSocket infrastructure is extended with two new channel types:
\begin{itemize}
    \item \textbf{Dispatch channel:} Broadcasts new delivery orders to empaneled drivers within the merchant's configured delivery radius. On native apps, incoming dispatches trigger push notifications via Capacitor's push-notifications plugin even when the app is backgrounded.
    \item \textbf{Tracking channel:} Streams driver GPS coordinates (sourced from Capacitor's geolocation plugin with high-accuracy mode) to both the merchant dashboard and the consumer's order tracking view.
\end{itemize}


\needspace{6\baselineskip}
\section{Risks \& Mitigation}

\begin{center}
\begin{tabularx}{\textwidth}{@{}l l X X@{}}
\toprule
\textbf{Risk} & \textbf{Likelihood} & \textbf{Impact} & \textbf{Mitigation} \\
\midrule
Insufficient driver supply at launch & Medium & Slow delivery times degrade consumer experience & Seed with merchant-employed delivery staff; phase launch by market (Pontevedra, FL first) \\
App Store review rejection & Low & Delays iOS launch by 1--2 weeks & Follow Apple Human Interface Guidelines; ensure native-feeling UX via Capacitor plugins; submit early to catch issues during Phase 4 \\
Brand-context bleed between apps & Low & Consumer sees driver UI or vice versa & Separate Capacitor projects with distinct bundle IDs; hostname-based layout isolation on web; automated E2E tests per touchpoint \\
AI-generated code quality gaps & Low & Subtle bugs in edge-case flows & Human review of all AI-generated code; comprehensive E2E testing on real devices; staged rollout to initial merchant cohort \\
Domain procurement delays & Low & Launch uses Surge URLs and native apps directly & Touchpoints are fully functional on Surge URLs and in native apps; vanity domains are cosmetic additions \\
\bottomrule
\end{tabularx}
\end{center}

\section{Success Criteria}

\begin{itemize}
    \item \textbf{Platform Availability:} Both apps published on App Store and Google Play Store within 30 days of project start.
    \item \textbf{Subscription Revenue:} At least 3 merchants with active \$50/month subscriptions within 30 days of launch (\$150 MRR).
    \item \textbf{Driver Fleet:} A minimum of 10 approved drivers (``Active'' status in Team panel) across all subscribing merchants within 60 days.
    \item \textbf{Order Volume:} 50+ delivery orders processed through BasaltDelivers within the first 90 days.
    \item \textbf{Uptime:} 99.5\% availability across all three surfaces (web, Android, iOS), measured via Cloudflare and app crash analytics.
    \item \textbf{Subscription Retention:} Less than 10\% monthly churn on BasaltDelivers subscriptions after the first 90 days.
\end{itemize}

\section{Requested Decision}

\begin{actionbox}
I am requesting approval to proceed with \textbf{Option A (Surge Touchpoints + Capacitor + Stripe Subscription)} at an estimated fixed budget of \textbf{\$154 in Year 1} (\$129/yr recurring) plus Stripe processing fees of 2.9\% + 30\textcent\ per subscription charge. Revenue target: \textbf{\$150+ MRR} within 30 days of launch. AI-accelerated engineering begins \textbf{immediately upon approval}, with all three surfaces (web, Android, iOS) targeted for production within \textbf{4 weeks}. Please confirm by \textbf{May 1, 2026}.
\end{actionbox}

% ============================================================================
% SIGNATURE
% ============================================================================

\vspace{2em}

\noindent\rule{5cm}{0.4pt}\\[4pt]
\noindent\textbf{\docAuthorName}\\
\noindent\textit{\docAuthorTitle}\\
\noindent BasaltHQ Inc.

\end{document}
